\exnumber{00}
\extitle{Projet 4 : À table !}
\exfiles{ex00/archive zip, nomme Name\_FirstName\_DeliverableName\_Date\_VersionNumber}
\exauthorize{Ceux nommés dans l'exercice}
\makeheaderfiles

Concevez une application de livraison de repas (interfaces et parcours utilisateur pour (1) les propriétaires de restaurants, (2) les clients, (3) les livreurs (4) les managers des livreurs, sur mobile et desktop).

DeliverNow est une entreprise qui aide les restaurants à livrer des plats à emporter, grâce à une équipe de runners.

Ils ont récemment reçu de mauvais commentaires de clients :

« Nos clients rencontrent actuellement des problèmes avec des délais de livraison estimés incorrects, la quantité de couverts et de condiments inutiles.

Nos runners se plaignent de la répartition des trajets entre eux.

Les managers n'arrivent pas à bien gérer leurs chauffeurs faute d'une vision suffisamment claire des runners.

Enfin, les restaurants se plaignent car ils ont du mal à gérer la préparation des plats sur place et à emporter.'

Ainsi, ils demandent à votre équipe de redesigner l'expérience complète :

\begin{itemize}
\item Pour les clients
\item Pour les restaurants
\item Pour les coureurs
\item Pour les managers de coureurs
\end{itemize}

Vous devez rencontrer votre client après 3 semaines pour une première présentation de votre travail, après avoir réalisé un prototype mid-fidelity testé pour chacun des 4 parcours utilisateur.

La présentation finale de votre travail aura lieu 2 semaines après la présentation intermédiaire, où vous montrerez un prototype haute fidélité testé de chacun des 4 parcours.

\textbf{Durée estimée :} 5 semaines

\textbf{Taille de l'équipe :} 3 personnes max

\textbf{Consignes :}

\begin{itemize}
\item Posez une stratégie UX et UI avant de foncer. Remettez le brief initial en question avec une recherche utilisateur.
\item Respectez les étapes : empathie, définir, énoncé du problème, idéation, prototype, test.
\item Dans ce contexte, 4 utilisateurs = 4 personas = 4 parcours utilisateurs = 4 prototypes
\item Vous devez interroger de vraies personnes pendant l'empathie, la recherche et les tests.
\item Dans vos prototypes, si nécessaire, montrez à votre public le parcours utilisateur de chaque persona lorsqu'il y a un ordre.
\item Dans les deux phases, testez vos prototypes avec au moins 5 personnes.
\item Faire les deux présentations comme si vous montriez votre travail à l'équipe DeliverNow. Assurez-vous qu'il y a au moins 5 personnes dans l'audience.
\item Préparez un formulaire de retour pour le public. Il contient au moins les questions listées plus bas.
\item Timing : 15 minutes pour la présentation et 5 minutes pour les Q\&R.
\item Vous pouvez vous servir des audits de présentation pour structurer vos présentations.
\item Déposez les documents sur Github, dans un zip nommé « ProjectTitle\_Name\_FirstName », et chaque livrable ainsi : « Name\_FirstName\_DeliverableName\_Date\_VersionNumber ». Exemple : « Doe\_John\_MidFiPrototype\_05242024\_V1 ».
\item Ajoutez un titre dans tous les documents écrits.
\item Soignez l'aspect global des livrables. Ils doivent être clairs, simples et lisibles. Vous pouvez regarder canva.com, sans surcharger le dessin de détails.
\item Amusez-vous !
\end{itemize}

\textbf{Outils :}

\begin{itemize}
\item Google Forms \url{https://docs.google.com/forms/u/0/}
\item Airtable \url{https://www.airtable.com/}
\item Typeform \url{https://www.typeform.com/}
\end{itemize}

\textbf{Pistes :}

\begin{itemize}
\item Compétences de présentation efficaces \url{https://academic.oup.com/femsle/article/364/24/fnx235/4587905}
\end{itemize}

Livrables intermédiaires :

\textbf{Citation :}

« Your mission in any presentation is to inform, educate, and inspire. » Robert Ballard

\medskip

\textbf{Phase 1 : audit de la présentation intermédiaire}

Vérifiez :

\begin{itemize}
\item Au moins 5 personnes assistent à la présentation
\item La présentation porte sur tout le processus UX, y compris :
\begin{itemize}
\item Le brief initial
\item Un planning, avec les phases UX, les outils choisis, le calendrier
\item Résultats de recherche utilisateur issus de la recherche qualitative et quantitative (persona, chiffres, etc.) pour les 4 utilisateurs
\item 4 énoncés de problème pour adapter le brief initial d'après les résultats de la recherche
\item Une introduction de chaque concept issu de l'idéation
\item Le prototype moyenne fidélité animé (une vidéo de démo) avec les 4 parcours utilisateurs
\item Les principaux retours des tests d'utilisabilité (5x4)
\item Les prochaines étapes du projet
\end{itemize}
\item La présentation ne dépasse pas 15 minutes
\item Les formulaires de retour ont été distribués au début de la présentation
\item Vous vous êtes amusés !
\end{itemize}

\medskip

\textbf{Phase 2 : audit de la présentation finale}

Vérifiez :

\begin{itemize}
\item Au moins 5 personnes assistent à la présentation
\item La présentation porte sur le processus UI, y compris :
\begin{itemize}
\item Un bref rappel de la présentation précédente
\item Un planning, avec les phases UI, les outils choisis, le calendrier
\item Une planche d'ambiance
\item Un styleguide (couleurs, marque, typographie, quelques éléments de la bibliothèque)
\item Le prototype haute fidélité animé (une vidéo de démo) avec les 4 parcours utilisateurs
\item Les retours principaux des tests d'utilisabilité et de désirabilité
\item Les prochaines étapes du projet
\end{itemize}
\item La présentation ne dépasse pas 15 minutes
\item Les formulaires de retour ont été distribués au début de la présentation
\item Vous vous êtes amusés !
\end{itemize}

\textbf{Questions à poser dans les formulaires de retour}

\begin{itemize}
\item Que retenez-vous de cette présentation ?
\item De 0 à 10, le discours était-il clair ?
\item Pourquoi ?
\item De 0 à 10, la présentation était-elle engageante ?
\item Pourquoi ?
\item De 0 à 10, la présentation était-elle cohérente ?
\item Pourquoi ?
\item De 0 à 10, les visuels portaient-ils ?
\item Pourquoi ?
\item De 0 à 10, le problème était-il clair ?
\item Pourquoi ?
\item De 0 à 10, le prototype répond-il au problème ?
\item Pourquoi ?
\item Quelle recommandation feriez-vous pour améliorer cette présentation ?
\item Pourquoi ?
\end{itemize}

